\section{Meta-RL 方法分类}

\subsection{Meta-RL 的典型任务设定}

Meta-RL 的经典例子是 maze navigation：先在一个新迷宫里收集少量试探性经验，再用这些经验解出这个迷宫。这里独有的难点是 exploration-exploitation trade-off 发生了变化: 早期交互不仅要赚取当前 reward，还要识别当前到底是哪一个任务。

\begin{figure}[H]
\centering
\includegraphics[width=0.9\textwidth]{slides-images/slide-10.jpg}
\caption{Meta-RL 例子：先在新迷宫里探索，再利用少量经验解决任务}
\end{figure}

\begin{figure}[H]
\centering
\includegraphics[width=0.9\textwidth]{slides-images/slide-11.jpg}
\caption{Meta-train / meta-test 的区别：训练时学会高效适应，测试时面对新任务快速上手}
\end{figure}

\subsection{更正式的问题表述}

这节课给了一个非常重要的 formal setup：Meta-RL 不是单个 MDP，而是一个\textbf{任务分布} $p(\mathcal{T})$。训练阶段从中抽取许多任务，学习一种适应机制；测试阶段在来自同一分布的新任务上，只给少量经验。

\begin{figure}[H]
\centering
\includegraphics[width=0.9\textwidth]{slides-images/slide-13.jpg}
\caption{Multi-task、transfer 与 meta-learning 的更正式定义}
\end{figure}

\begin{importantbox}{Meta-RL 的输入输出}
以 episodic 版本为例：
\begin{itemize}
    \item 输入：来自某个任务的少量 rollout / experience $D_{\text{train}}$
    \item 输出：一个根据 $D_{\text{train}}$ 调整后的策略 $\pi(a\mid s, D_{\text{train}})$
\end{itemize}
也就是说，传统多任务 RL 中显式给出的 task ID，在 Meta-RL 里被一小段交互经验替代了。
\end{importantbox}

\begin{figure}[H]
\centering
\includegraphics[width=0.9\textwidth]{slides-images/slide-14.jpg}
\caption{Meta-RL 的输入输出：经验上下文代替显式 task descriptor}
\end{figure}

\subsection{方法 1：基于上下文的方法（Context-Based）}

核心思想是将新任务的经验作为\textbf{上下文}输入给策略网络，让网络根据上下文推断任务并做出决策。

\begin{importantbox}{上下文策略}
$$\pi_\theta(a | s, c)$$
其中 $c = \{(s_1, a_1, r_1, s_1'), \ldots, (s_k, a_k, r_k, s_k')\}$ 是在新任务上收集的少量经验。策略网络需要从这些经验中推断出任务的特性。

常用架构：
\begin{itemize}
    \item 将上下文通过 RNN/LSTM 编码为隐状态
    \item 将上下文通过 Transformer 的注意力机制编码
    \item 将上下文通过集合编码器（permutation-invariant）编码
\end{itemize}
\end{importantbox}

\begin{knowledgebox}{RL\textsuperscript{2} (Learning to Reinforcement Learn)}
RL$^2$ 是最早的 context-based meta-RL 方法之一。它将整个 RL 过程视为一个更大的 POMDP，其中 RNN 的隐状态编码了关于当前任务的"信念"。外层 RL 算法训练 RNN 使其能在内层快速适应新任务。
\end{knowledgebox}

\subsection{Black-box Meta-RL 的训练循环}

这门课重点讲的是 black-box meta-RL：不显式写出某个贝叶斯更新公式，而是直接用一个带记忆的神经网络，把 ``如何根据经验适应'' 这件事学进参数里。

\begin{figure}[H]
\centering
\includegraphics[width=0.9\textwidth]{slides-images/slide-17.jpg}
\caption{Black-box Meta-RL 的基本算法：按任务收集多轮 episode，再统一更新带记忆策略}
\end{figure}

\begin{knowledgebox}{它和普通 recurrent policy 有什么不同}
表面上看，black-box meta-RL 很像 ``带 RNN 的 RL''，但课程特别强调了两个额外条件：
\begin{enumerate}
    \item reward 通常作为输入的一部分进入网络；
    \item 隐状态需要在同一任务的多个 episode 之间持续保留。
\end{enumerate}
因此网络记住的不只是单条轨迹，而是 ``我已经在这个任务上观察到了什么''。
\end{knowledgebox}

\subsection{常见架构与优化器}

black-box meta-RL 的实现空间很大：可以用 RNN、attention、Transformer，也可以配 TRPO/A3C 或 SAC 等不同 outer-loop RL 优化器。

\begin{figure}[H]
\centering
\includegraphics[width=0.9\textwidth]{slides-images/slide-19.jpg}
\caption{Black-box Meta-RL 的架构与优化器选择：RNN、attention、off-policy context variable 等}
\end{figure}

\subsection{方法 2：基于梯度的方法（Gradient-Based）}

\begin{importantbox}{MAML（Model-Agnostic Meta-Learning）}
MAML 的核心思想是学习一组好的初始化参数 $\theta$，使得对任何新任务，只需少量梯度步即可适应。

\textbf{内层更新}（任务适应）：
$$\theta_i' = \theta - \alpha \nabla_\theta \mathcal{L}_{\mathcal{T}_i}(\theta)$$

\textbf{外层更新}（元学习）：
$$\theta \leftarrow \theta - \beta \nabla_\theta \sum_i \mathcal{L}_{\mathcal{T}_i}(\theta_i')$$

关键：外层优化的目标是让\textbf{适应后}的参数 $\theta_i'$ 在任务 $\mathcal{T}_i$ 上表现好，而非原始参数 $\theta$。
\end{importantbox}

\begin{warningbox}{MAML 的计算开销}
MAML 需要对"梯度"求梯度（二阶导数），计算开销较大。虽然可以用一阶近似（如 FOMAML、Reptile）来降低开销，但这会牺牲一定的精度。在 RL 中，由于策略梯度本身方差就很高，MAML 的训练稳定性是一个挑战。
\end{warningbox}

\subsection{两类方法的比较}

\begin{table}[H]
\centering
\begin{tabular}{lcc}
\toprule
& \textbf{Context-Based} & \textbf{Gradient-Based} \\
\midrule
适应方式 & 前向传播（推断） & 梯度下降 \\
适应速度 & 极快（一次前向） & 需要若干梯度步 \\
表达能力 & 受限于网络容量 & 理论上可精确适应 \\
训练难度 & 相对容易 & 需要二阶优化 \\
典型方法 & RL$^2$, PEARL & MAML, ProMP \\
\bottomrule
\end{tabular}
\caption{两类 Meta-RL 方法的比较}
\end{table}

\subsection{Black-box Meta-RL 的案例}

课程给了两个很有代表性的案例。第一个是视觉迷宫导航：模型在大量小迷宫上 meta-train，测试时要在看不见过的迷宫里快速探索并找到出口。第二个案例更贴近当前大模型语境：把不同数学题看作不同任务，优化 test-time compute，让模型在固定 token 预算下更高效地尝试、验证并选择策略。

\begin{figure}[H]
\centering
\includegraphics[width=0.9\textwidth]{slides-images/slide-20.jpg}
\caption{Meta-RL 案例 1：视觉迷宫导航，检验模型能否在新迷宫里快速形成策略}
\end{figure}

\begin{figure}[H]
\centering
\includegraphics[width=0.9\textwidth]{slides-images/slide-22.jpg}
\caption{Meta-RL 案例 2：把不同数学题视为任务，优化 test-time compute}
\end{figure}

\subsection{Meta-RL 与多任务策略的关系}

讲者专门用一页 slides 解释了一个容易混淆的点：多任务策略依赖显式 task ID $z_i$，而 Meta-RL 把\textbf{经验本身}当成 task identifier。换句话说，多任务策略是 ``我告诉你现在是哪一类任务''，而 Meta-RL 是 ``你自己从交互里推断出来''。

\begin{figure}[H]
\centering
\includegraphics[width=0.9\textwidth]{slides-images/slide-24.jpg}
\caption{当 task ID 不在状态里时，策略必须从经验中推断当前所处的 MDP}
\end{figure}

\subsection{本章小结}

Meta-RL 有两大技术路线：context-based 方法通过编码经验实现隐式适应，gradient-based 方法通过学习好的初始化实现快速微调。

